\chapter{相关理论与技术}\label{chap:guide}

\section{检索增强生成技术综述}

\subsection{RAG基本范式}

检索增强生成（Retrieval-Augmented Generation, RAG）是一种结合了参数化记忆与非参数化记忆的混合框架~\citep{lewis2020rag,guu2020realm,gao2024retrievalaugmented}。其核心思想是：在生成回答之前，首先根据输入的查询从海量知识库中获取相关信息，并将其与原始输入结合，共同作为语言模型的上下文输入。

形式化地定义，给定一个用户查询 $q$，RAG 系统的处理流程可以表示为以下三个核心步骤：

（1）检索阶段： 系统利用检索器 $f_{ret}$ 从外部知识库 $\mathcal{D} = \{d_1, d_2, \dots, d_n\}$ 中筛选出最相关的 $k$ 个文档片段 $\mathcal{D}_{top-k}$。其公式如下：
\begin{equation}
    \mathcal{D}_{top-k} = \text{top-k}(\{ \text{sim}(q, d_i) \mid d_i \in \mathcal{D} \})
\end{equation}
其中，$\text{sim}(\cdot, \cdot)$ 是衡量查询与文档之间语义相关性的相似度函数。

（2）增强阶段：将检索到的文档 $\mathcal{D}_{top-k}$ 按照特定模板与原始查询 $q$ 进行拼接，构造出一个增强后的上下文 $\mathcal{P}$（即 Prompt）。其公式如下：
\begin{equation}
    \mathcal{P} = \text{Template}(q, \mathcal{D}_{top-k})
\end{equation}

（3）生成阶段： 语言模型 $P_{LM}$ 基于增强后的上下文 $\mathcal{P}$ 生成最终的回复 $y$。公式如下所示：
\begin{equation}
    y \sim P_{LM}(y \mid \mathcal{P}) = \prod_{t=1}^T P_{LM}(y_t \mid \mathcal{P}, y_{<t})
\end{equation}

在实际应用中，现代 RAG 系统多采用密集向量检索（Dense Retrieval）技术~\citep{karpukhin-etal-2020-dense,izacard2022contriever,xiao2023cpack}。通过预训练的编码器（Encoder）将查询和文档映射至高维向量空间 $\mathbb{R}^d$，使得语义层面的匹配能够通过计算向量间的距离来实现。

\subsection{RAG中检索模块的关键作用}

在 RAG 框架中，检索模块不仅是信息的来源，更是决定整个系统性能上限的关键组件。其重要性主要体现在以下三个方面：

（1）缓解事实性幻觉。大语言模型容易在生成过程中产生与事实不符的内容~\citep{xu2025hallucination,shuster2021retrieval}。检索模块通过提供来自可信源（如 Wikipedia、企业私有数据库）的外部证据，将生成过程约束在已知的事实范围内。其公式如下：
\begin{equation}
    P(y \mid q, \mathcal{D}_{top-k}) > P(y \mid q)
\end{equation}
当 $\mathcal{D}_{top-k}$ 包含正确答案时，模型输出真实事实的后验概率显著提升。

（2）突破知识时效性限制。传统 LLM 的知识被锁定在模型参数 $\theta$ 中，更新成本极高。检索模块通过解耦知识与推理，使得系统只需更新外部索引（Index）即可获取最新信息，而无需对模型进行微调。这使得 RAG 系统在处理突发事件或实时信息时具有天然的优势。

（3）提升长尾与专业领域表现。研究表明，LLM 对训练语料中低频出现的“长尾实体”记忆效果较差~\citep{kandpal2023largelanguagemodelsstruggle}。检索模块通过显式地将长尾知识引入上下文，弥补了模型参数记忆的不足。在医疗、法律等专业领域，这种“开卷考试”式的范式能大幅提高系统在特定垂直领域的专业性与证据可溯源性。

综上所述，检索模块的召回率和精确率直接决定了 RAG 系统能否有效利用外部知识。若检索模块失效，生成模型将面临“无米之炊”或“错误引导”的困境，这正是本文研究多轮对话检索评估基准的根本动因。



\section{对话检索基准数据集}

对话检索（Conversational Information Retrieval, CIR）的评估效能取决于基准数据集对复杂交互逻辑的还原程度。在 RAG 技术演进的背景下，数据集构造模式经历了从受限环境下的语义理解到大规模开放域检索，再到当前拓扑驱动的自动化生成三个阶段。本节将深入探讨各阶段代表性数据集的构造机理及其在处理多轮上下文依赖时的内在局限。

\subsection{人工标注数据集}

早期对话检索基准的构建逻辑立足于对人类真实信息寻求行为的“认知还原”。其核心目标是在受限语料环境下，捕捉自然语言中的指代、省略以及话题细微偏移等语义现象。

\begin{itemize}
    \item \textbf{QuAC \citep{choi-etal-2018-quac}：} 该数据集通过“师生博弈”模式引入了信息不对称机制。提问者在不可见背景文档的前提下进行试探性提问，而回答者必须根据上下文判断回复内容或引导对话。QuAC 的重要贡献在于提出了“不可回答”标签与“对话延续性”反馈，强迫模型建模对话的流状态。
    \item \textbf{CoQA \citep{reddy-etal-2019-coqa}：} 相比 QuAC，CoQA 强调了回答的自由度与多领域泛化。其标注逻辑允许非片段抽取的回复，且数据涵盖了从文学到新闻的多样化语料。其分析显示，模型在长程对话中极易丢失对早期实体的感知，揭示了神经网络在处理长上下文时的“中间遗忘”特性。
    \item \textbf{Doc2Dial \citep{feng-etal-2020-doc2dial}：} 该数据集代表了任务导向型构造逻辑。其深度在于将对话流与文档的层级拓扑进行显式硬映射。标注过程模拟了用户在专业领域（如政务或法律）中按照文档逻辑链条逐步深入的导航路径，定义了对话行为受制于知识结构的特征。
\end{itemize}

然而，人工标注模式在评估 RAG 系统时面临严重的“认知边界”与“标注稀疏性”陷阱。标注者通常在“读后提问”的局部视野下工作，仅针对当前阅读的文档进行标注。这种构造机理导致了一个根本性的矛盾：在大规模语料库环境下，大量客观存在且同样能够支持答案的证据文档由于未被标注者预阅，而被系统判定为负样本。这种评估偏置使得模型在开放域场景下的召回率表现被系统性低估，限制了对检索器真实性能的客观度量。


\subsection{半自动与混合标注数据集}

随着研究重心向大规模开放域转移，研究者开始利用已有检索模型构建初始候选集，再由人工进行二次标注或查询重写。

\begin{itemize}
    \item \textbf{TREC CAsT \citep{Dalton2020TRECC2}：} 该基准利用 池化判定模式 解决了千万级语料库的评估公信力问题。它汇聚了多个参赛系统的检索结果，并由专家进行多级相关性判别。这种模式的意义在于建立了不依赖于单一模型偏见的可重用测试集，是衡量复杂搜索任务鲁棒性的核心基准。
    \item \textbf{QReCC \citep{anantha-etal-2021-qrecc}：} 其构造核心是 “去语境化”查询改写 。通过人工将模糊查询转化为语义自洽的独立查询，QReCC 将端到端对话检索任务解构为“改写+标准检索”的级联流程。这种模式支持对 RAG 系统中检索模块与改写模块的解耦评估，极大地便利了系统的错误归因。
\end{itemize}

尽管混合模式显著提升了检索空间的覆盖度，但其本质是“广度优先”的工程妥协。池化方法天然地对“多数派”算法存在依赖，对具备独特表述理解机制的创新检索器可能产生判别偏置。此外，QReCC 式的解耦逻辑虽然降低了工程难度，却在某种程度上割裂了对话的整体性。由于忽略了交互过程中意图的动态演化与反馈循环，该模式构建的基准往往呈现出一种“静态化”特征，难以模拟真实 RAG 系统在面对多轮追问时的动态不确定性。


\subsection{全自动合成数据集}

针对人工标注的高成本与低扩展性问题，利用大语言模型自动化模拟对话路径并生成大规模基准成为当前趋势。CORAL~\citep{Cheng2024CORALBM} 是该模式的早期代表，其核心机制为：
\begin{itemize}
    \item \textbf{层级标题树遍历：} CORAL 利用维基百科的目录层级结构（标题-子标题-正文），通过 LLM 将章节标题改写为用户查询，正文提取为答案，引用文献作为黄金文档。
    \item \textbf{多维归因标注：} 尝试建立“检索-生成-引用”的联合标注逻辑，要求模型指明生成内容对证据文档的依赖关系。
\end{itemize}

尽管 CORAL 在自动化规模上取得突破，但其质量缺陷已被实证研究揭露。CORAL 数据集存在严重的系统性错位：生成的对话内容（如天文学问题）与其标注的支持文档（如 DOI 标识符系统说明）在语义上毫无关联。这种“盲视生成”源于其过度依赖静态文档结构，缺乏全局语义验证机制。由于该问题无法通过简单规则修复，CORAL 团队已主动撤回原版数据集。这一案例深刻表明：脱离全局感知与意图驱动的自动化合成，仅能产出规模庞大但信噪比极低的“数据废料”，无法胜任严谨的检索能力评估。



\section{基于LLM的数据评估与生成方法}
随着大语言模型（LLMs）推理与指令遵循能力的成熟，其在基准数据构建与质量评估中正从辅助工具转变为核心驱动引擎。这不仅改变了传统依赖昂贵人工标注的作坊式生产模式，更为突破“认知边界”与“标注稀疏性”提供了全新的技术路径。本小节将系统梳理当前基于LLM的自动化评估与多智能体生成两大关键技术范式。
\subsection{LLM-as-a-Judge评估范式}
在传统信息检索与自然语言处理评估中，精确的人类专家判定是衡量系统性能的“黄金标准”，但其高昂的时间与经济成本严重制约了迭代速度。LLM-as-a-Judge范式~\citep{zheng2023judging,liu2023gpteval}的兴起，旨在利用先进大语言模型（如GPT-4~\citep{openai2023gpt4}、DeepSeek-V3~\citep{deepseekai2024deepseekv3}）的内在推理能力，模拟人类对文本质量的综合评判。

该范式的核心运作机理是将评估任务重构为特定的提示工程问题。给定待评估的对象（如检索结果的相关性、生成答案的准确性），系统构造包含详尽评分准则与参考证据的提示词，引导LLM执行成对比较或绝对评分。为了提高判别精度，研究者广泛引入了思维链（Chain-of-Thought, CoT）技术~\citep{wei2022chain}，强制模型先输出推理逻辑再给出最终结论，有效减少了直觉性误判。

在对话检索基准的语境下，LLM-as-a-Judge展现了显著的应用价值。相较于早期仅依赖字符串重叠度（如 BLEU~\citep{papineni2002bleu}、ROUGE~\citep{lin2004rouge}、METEOR~\citep{banerjee2005meteor}）或简单二元匹配的指标，LLM能够结合对话上下文，灵活处理语义等价但表述迥异的答案，并对证据支持度做出深层次判断。这为解决多轮对话中的指代消解与话题隐含依赖评估提供了可行的自动化解决方案。
然而，将该范式直接应用于严苛的检索基准质量评估时，仍暴露出三处亟待解决的系统性局限：

（1）自偏好偏差~\citep{panickssery2024llm}：LLM评估者倾向于给由自身或同族模型生成的内容打更高分数。当基准数据恰好由同一家族模型合成时，这种偏差会导致评估结果虚假膨胀，形成“模型内卷”的评测泡沫。

（2） 知识幻象干扰：大语言模型自身的参数化知识可能与外部提供的上下文发生冲突。当评估涉及冷门事实或最新资讯时，模型固有的过时知识或幻觉可能导致其无视给定的证据文档，转而依赖错误的内部记忆进行评判，造成评估信度崩塌。

（3）指令敏感性：评分结果对提示词的措辞、评分量表的粒度（如5分制vs10分制）以及输出格式极其敏感。细微的提示变动可能引发评分分布的剧烈漂移，使得跨研究的评估结果难以复现与对齐。

为克服上述缺陷，本文在 CRB-Eval 中采用多模型集成策略与逐点评分设计。通过在多文档判别任务中随机打乱文档位置，彻底消除位置偏差；通过聚合异构模型的评分分布，有效稀释单一模型的自偏好倾向，使跨家族裁判一致性（以 Fleiss $\kappa$ 度量）稳定提升至 \TBD{提升后$\kappa$下限} 以上。


\subsection{多智能体对话生成}
面对多轮对话检索中复杂的意图演化与上下文缠绕，传统的单步文本改写难以胜任。多智能体对话生成通过模拟真实的信息交换生态，为大跨度、高保真的对话流合成开辟了新路径。

该方法摒弃了单一模型单向输出的模式，转而设立多个具备差异化角色的智能体~\citep{wang2024survey}。典型配置包括：用户模拟器，负责根据潜在知识需求发起追问；系统模拟器，模仿RAG系统执行检索与生成应答；以及可能的评判员在后台校验事实一致性。各智能体在共享对话历史的状态空间下交替轮转，通过多轮对抗与合作，自发涌现出包含澄清、纠错与话题漂移的完整会话轨迹。Park等人的开创性工作展示了多智能体环境在模拟人类社会互动方面的惊人潜力~\citep{Park2023GenerativeAI}，也为合成高质量对话评估数据提供了理论背书。

在基准构建实践中，此类架构的优势在于能够实现闭环可控性。研究者可以通过设定用户画像与目标知识图谱，精准诱导对话流向特定的考察难点（如多跳推理、专有名词指代）。相比于早期CORAL~\citep{coral2025}依赖静态文档拓扑的生成方式，多智能体仿真更贴近人类思维的跳跃性与容错性，生成的查询往往具备更强的口语化扰动与突发性话题转折，对检索模型的鲁棒性构成更大挑战。

尽管如此，当前的多智能体合成方案在服务于工业级检索评估时仍存在显著短板。首当其冲的是模拟器的真实性鸿沟。现有模拟器多基于强指令调整模型构建，其产生的对话往往呈现过度礼貌、逻辑过度严密、极少犯错的“教科书式”特征。这种缺乏语言噪音与混乱意图流的“温室对话”，导致训练出的检索模型在面对真实用户的模糊、破碎输入时表现骤降。此外，多数现有框架缺乏对生成内容的细粒度溯源约束，智能体可能在未严格绑定外部证据的情况下流畅续写，加剧了合成数据的本体性幻觉风险，削弱了基准的可信度。

鉴于此，本文在第5章提出CRB-Pipeline流水线，旨在吸收多智能体分工协作的效率优势，同时通过引入混合质量过滤与角色间相互制衡机制，深度压制上述偏差，以产出兼具自然度与证据严谨性的高质量基准数据。



\section{文本聚类与语义路径构建}

在自动化对话检索基准的生成过程中，如何从海量非结构化语料库中提取具备逻辑关联的知识片段，并将其组织为符合人类认知逻辑的会话序列，是决定基准质量的关键。本节将从文本语义表示、聚类分析以及语义路径构建三个维度综述相关理论与技术。

\subsection{文本语义表示与向量空间模型}

文本聚类的基础在于对文本信息的数学化表征。传统的词袋模型（BoW）与 TF-IDF 算法虽然在稀疏检索中表现稳健，但难以捕捉深层的语义关联。随着深度学习的发展，基于预训练语言模型的分布式表示已成为主流。

给定文本分块 $c_i$，通过编码器 $E$（如 BERT~\citep{devlin2019bert}）映射至高维嵌入空间 $\mathbb{R}^d$。其公式如下：
\begin{equation}
    \mathbf{v}_i = E(c_i)
\end{equation}

在此向量空间中，文本间的语义相似度通常由欧氏距离或余弦相似度衡量。为了提升聚类的计算效率并消除高维空间的“维度灾难”，近似最近邻搜索（ANN）与流形降维（如 UMAP、t-SNE）常被用于辅助文本分布的拓扑分析。这些表征技术为后续在海量知识库中通过聚类识别“知识主题簇”奠定了物理基础。

\subsection{文本聚类算法的演进与局限}

文本聚类旨在将语义相近的文档片段归并为同一主题集合，从而为对话生成的上下文提供候选池。

（1）基于质心的聚类：以 K-Means 为代表。其核心逻辑是通过最小化簇内平方误差来识别中心点。然而，该算法要求预验指定聚类数目 $k$，且在处理边界模糊、形状不规则的文本流形时效果较差。

（2）基于密度的聚类：以 DBSCAN 为代表。该类算法通过识别高密度区域来发现任意形状的簇，并能有效识别离群点（即噪声）。其优势在于无需预设簇数，但对距离阈值（$\epsilon$）极为敏感，在全局数据密度分布不均时难以获得理想的划分。

（3）层次聚类： 通过构建树状拓扑反映文本间的包含关系。虽然其解释性强，但在面对工业级规模的知识库时，其 $O(n^2)$ 或 $O(n^3)$ 的时间复杂度使其难以在生产环境下落地。

在对话检索基准合成的语料预处理阶段，现有算法往往难以平衡“簇内语义高内聚”与“簇间逻辑强关联”的双重需求。K-Means 生成的簇大小不可控，易超出 LLM 上下文窗口；DBSCAN 对密度阈值敏感；层次聚类时间复杂度高达 $O(n^3)$。更关键的是，基于相似度阈值的传统聚类会导致不同簇之间出现大量文档重叠，使 LLM 在训练与评估中重复接触相同内容，引入严重的评估偏差。



\section{本章小结}

本章对本文研究工作所涉及的核心理论与关键技术进行了系统性的梳理与综述。首先，阐述了检索增强生成（RAG）的基本范式，明确了检索模块在缓解事实性幻觉、提升知识时效性方面的决定性作用，从而确立了对对话检索模块进行科学评估的必要性。其次，通过对现有对话检索基准数据集的横向对比，揭示了人工标注范式在认知边界上的局限、混合标注范式在检索规模上的妥协，以及全自动合成范式在自然度与证据严谨性方面的不足。随后，重点探讨了基于大语言模型（LLM）的自动化评估与生成方法，分析了 LLM-as-a-Judge 范式下的各类系统偏见，以及多智能体协作生成在模拟复杂交互意图时的潜力与挑战。最后，论述了文本语义表示、聚类算法及语义路径构建的图论模型，说明了将对话流生成抽象为基于语义流形的路径优化问题的可行性。

通过上述理论回顾可以发现，尽管 RAG 技术与多轮对话系统已取得显著进展，但如何构建一个既能突破人工标注认知局限、又能克服自动化生成僵化缺陷的高质量、规模化评估基准，并提供标准化的审计手段，依然是当前领域亟待填补的空白。本章的分析为后文提出 CRB-Suite 统一框架，包括 CRB-Eval 审计方法、CRB-Pipeline 合成流水线以及 CRB-Bench 基准数据集，奠定了扎实的理论基础与技术依据。

